
\assignq{1}{2}{Beam splitter on Fock and coherent inputs}{1}
The input/output relation is the real rotation
\[
\begin{pmatrix}c\\d\end{pmatrix}
=\begin{pmatrix}\tau&-\varrho\\\varrho&\tau\end{pmatrix}
\begin{pmatrix}a\\b\end{pmatrix}.
\]
Unitarity means \(\tau^2+\varrho^2=1\).
The coefficients appearing in the \emph{matrix}
are field-amplitude coefficients: the corresponding
intensity fractions are \(\tau^2,\varrho^2\).

\subsection*{(a) The binomial transformation}
The inverse mode relation gives
\(a^\dagger=\tau c^\dagger+\varrho d^\dagger\).
Write the number state with creation operators:
\[
\ket{n,0}_{ab}=\frac{(a^\dagger)^n}{\sqrt{n!}}\ket0.
\]
Expand the commuting creation operators:
\[
\frac1{\sqrt{n!}}\sum_{k=0}^n\binom nk
\tau^k\varrho^{n-k}
(c^\dagger)^k(d^\dagger)^{n-k}\ket0.
\]
Because the normalised number states satisfy
\((c^\dagger)^k(d^\dagger)^{n-k}\ket0
=\sqrt{k!(n-k)!}\ket{k,n-k}\), this becomes
\[
\boxed{\ket{n,0}_{ab}\longrightarrow
\sum_{k=0}^n
\sqrt{\frac{n!}{k!(n-k)!}}\,
\tau^k\varrho^{n-k}\ket{k,n-k}_{cd}}.
\]
Normalisation follows from
\(\sum_k\binom nk\tau^{2k}\varrho^{2(n-k)}
=(\tau^2+\varrho^2)^n=1\).

\subsection*{(b) Coherent input plus vacuum}
For \(c=(a+b)/\sqrt2,\ d=(a-b)/\sqrt2\),
write the input as
\[
\ket{\alpha,0}=e^{-|\alpha|^2/2}
\exp(\alpha a^\dagger)\ket0.
\]
The inverse relation is
\(a^\dagger=(c^\dagger+d^\dagger)/\sqrt2\).
The two output creation operators commute, so
the exponential factors:
\[
\begin{aligned}
\ket{\Psi_{\rm out}}
&=e^{-|\alpha|^2/2}
e^{(\alpha/\sqrt2)c^\dagger}
e^{(\alpha/\sqrt2)d^\dagger}\ket0\\
&=\boxed{\ket{\alpha/\sqrt2}_c\otimes
\ket{\alpha/\sqrt2}_d}.
\end{aligned}
\]
Phase conventions can change the sign of one
amplitude without changing output count statistics.

\subsection*{(c) Photodetection}
The \(c\)-mode coherent amplitude is
\(\beta=\alpha/\sqrt2\).
Thus
\[
\boxed{P_c(n)=
e^{-|\alpha|^2/2}\frac{(|\alpha|^2/2)^n}{n!}},
\qquad
\boxed{\langle n_c\rangle=\operatorname{Var}(n_c)
=\frac{|\alpha|^2}{2}}.
\]
The ideal detector observes \emph{Poissonian} counts.
\methodnote{Normalised Fock coefficients need the
square roots of factorials; keep amplitude and intensity
transmission coefficients distinct.}
